% ============================================================================
% INFORMACJA O WYKORZYSTANIU NARZĘDZI AI
% ============================================================================
% Ta część realizuje wzór z załącznika nr 2A do Zasad korzystania z narzędzi
% AI przez studentów i doktorantów Politechniki Opolskiej.
%
% Wybierz dokładnie jeden wariant:
%   1) pozostaw treść poniżej, jeżeli AI wpłynęła na treść, strukturę,
%      rozwiązanie, kod, obliczenia, analizę, interpretację, wnioski lub inny
%      element materiału przekazywanego do oceny;
%   2) zastąp ją wariantem umieszczonym na końcu pliku, jeżeli nie korzystano
%      z AI w zakresie wymagającym ujawnienia.
%
% Zwykła autokorekta, sprawdzanie pisowni i inne funkcje techniczne, które nie
% wpływają na treść merytoryczną, co do zasady nie wymagają ujawnienia, chyba że
% prowadzący określi inaczej. Zakres użycia AI musi być zgodny z zasadami
% ustalonymi przez prowadzącego. Na jego żądanie trzeba przedstawić bardziej
% szczegółowe informacje, np. przykładowe polecenia, wyniki lub materiały
% robocze.
% ============================================================================

\chapter{Wykorzystanie narzędzi sztucznej inteligencji}
\label{chap:wykorzystanie-ai}

% === WARIANT 1: KORZYSTANO Z AI W ZAKRESIE WYMAGAJĄCYM UJAWNIENIA ===
% Dostosuj formy „korzystałem/am” i „przeze mnie” albo zastosuj formy mnogie,
% jeżeli projekt przygotował zespół.

W przygotowaniu niniejszego projektu korzystaliśmy z następujących narzędzi AI:

\noindent\textbf{Nazwa narzędzia lub narzędzi:}\par
Nazwa narzędzia, model lub wersja.

\medskip
\noindent\textbf{Cel i zakres wykorzystania:}\par
Opisz cel i zakres użycia, np. tworzenie pomysłów, analiza kodu, korekta
językowa albo proponowanie przypadków testowych.

\medskip
\noindent\textbf{Elementy przygotowane lub zmienione z wykorzystaniem AI:}\par
Wskaż konkretne rozdziały, fragmenty tekstu, kodu, obliczeń, analiz,
wizualizacji lub inne elementy projektu.

Wyniki uzyskane przy użyciu AI zostały przez nas sprawdzone, a za ostateczną
treść i poprawność przygotowanego materiału ponosimy odpowiedzialność.

% Nie wolno wprowadzać do niezatwierdzonych przez Uczelnię narzędzi AI m.in.
% danych osobowych, informacji poufnych i materiałów niepublicznych,
% nieopublikowanych danych i wyników badań, materiałów egzaminacyjnych, prac
% innych studentów, haseł, tokenów ani kluczy dostępu. Należy także przestrzegać
% praw autorskich, licencji, prawa do wizerunku i warunków użycia narzędzi.

% === WARIANT 2: NIE KORZYSTANO Z AI W ZAKRESIE WYMAGAJĄCYM UJAWNIENIA ===
% Aby użyć wariantu 2, usuń lub zakomentuj cały wariant 1 powyżej, a następnie
% odkomentuj poniższy akapit. Dostosuj formę osobową do autora albo zespołu.
%
% Oświadczamy, że podczas przygotowywania niniejszego projektu nie
% korzystaliśmy z narzędzi AI w sposób, który wpłynął na treść, strukturę,
% rozwiązanie, kod, obliczenia, analizę, interpretację, wnioski lub inne
% elementy materiału przekazywanego do oceny.
